\section{Микромир}
\label{sec:macro-open}
\label{sec:microworld}

Внутри вещества непрерывные уравнения, только что записанные, ещё держат механику и электродинамику прибора.
Они перестают держать сам опыт, как только отсчёт становится единичным.

Ансамбль подготовленных одинаково атомов даёт частоты исходов.
Один запуск прибора даёт один исход.
Уравнения поля содержат частоты и не содержат границу между системой, прибором и записью результата.
Сохранение того, что в ансамбле выглядит как полная вероятность, здесь ещё не объясняет одиночный отсчёт.

Числа~$\speedC$, $\planckHbar$, $\constGrav$, $\chargeE$ и массы фермионов измерены с большой точностью.
Безразмерное число, которое из них собирается, --- постоянная тонкой структуры
\begin{equation}
  \alphaFs = \frac{\chargeE^2}{4\symPi\constEpsZero\planckHbar \speedC}.
  \label{eq:alpha-fs}
\end{equation}
\begin{align}
  [\alphaFs]
  &= \left[\frac{\chargeE^2}{4\symPi\constEpsZero\planckHbar \speedC}\right]
  = \frac{\mathrm{Кл}^2}{\bigl(\NotationUnitFaradPerMetre\bigr)\,(\mathrm{Дж\cdot с})\,\bigl(\NotationUnitMetrePerSecond\bigr)}
  \notag\\
  &= \frac{\mathrm{Кл}^2}{\bigl(\NotationUnitFarad\bigr)\,(\mathrm{В\cdot Кл})}
  = 1.
  \label{eq:alpha-fs-dim}
\end{align}
Coupling безразмерен: фарады, джоули и скорость в знаменателе сокращаются до единицы.
Она стоит в законе свободным множителем:
измерение есть, вывода из \eqref{eq:maxwell-hom}--\eqref{eq:newton} нет.
Пока закон претендует только описать прибор, такой множитель терпим.
Как только тот же закон считается фундаментальным, вписанное число означает, что измерение не выведено.

\begin{remark}[Формула \eqref{eq:alpha-fs} --- мост, не онтология]
\label{rem:alpha-fs-bridge}
Запись $\alphaFs=\chargeE^2/(4\symPi\constEpsZero\planckHbar \speedC)$ удобна на~$\layerMacro$:
безразмерное число \emph{собрано} из измеренных размерных величин.
На~$\layerMicro$ тот же объект --- coupling фазы и вакуума и отношение квантов силы $\ladderForce$
(\cref{ch:alpha}, \cref{def:alpha-meaning}); $\chargeE$ не входит в его \emph{определение}.
\end{remark}

Чёрная дыра делает ограниченность плотности видимой числом.

\begin{definition}[Битовый бюджет]
\label{def:bit-budget}
Пусть у системы задана термодинамическая энтропия~$\entropyS$.
\textbf{Битовым бюджетом} $B$ говорят о величине
\[
  B := \frac{\entropyS}{\constBoltzmann\ln 2},
\]
то есть об энтропии, пересчитанной в биты (безразмерно).
\end{definition}

\begin{definition}[Число различимых микросостояний]
\label{def:distinguishable-microstates}
Пусть у изолированной системы энтропия~$\entropyS$.
Под \textbf{числом различимых микросостояний} $\opNormN$ понимают величину,
удовлетворяющую $\entropyS=\constBoltzmann\ln\opNormN$:
это логарифм числа микросостояний, согласованных с макроописанием.
\end{definition}

\begin{definition}[Предел Бекенштейна]
\label{def:bekenstein-bound}
Пусть причинно связанная область вписана в шар радиуса~$R$,
имеет суммарную энергию~$\energyE$ (материя и поля),
а $A$ --- площадь её минимальной причинной границы.
Тогда \textbf{предел Бекенштейна} для такой области --- это система верхних оценок
на энтропию~$\entropyS$, на $\ln\opNormN$ (\cref{def:distinguishable-microstates})
и на битовый бюджет~$B$ (\cref{def:bit-budget});
конкретные неравенства приведены в \cref{prop:bekenstein-inequalities}.
\end{definition}

\begin{proposition}[Неравенства Бекенштейна]
\label{prop:bekenstein-inequalities}
Планковская площадь
\[
  \lengthPlanck^2=\frac{\planckHbar\,\constGrav}{\speedC^3}.
\]
Для области из \cref{def:bekenstein-bound}:
\begin{align}
  \entropyS &\le \frac{2\symPi\,\constBoltzmann\, R\,\energyE}{\planckHbar\,\speedC},
  \label{eq:bekenstein-S}\\
  B &\le \frac{2\symPi\, R\,\energyE}{\planckHbar\,\speedC\,\ln 2},
  \label{eq:bekenstein-bits}\\
  \ln\opNormN &\le \frac{\entropyS}{\constBoltzmann}
  \le \frac{A}{4\lengthPlanck^2}.
  \label{eq:state-count}
\end{align}
На горизонте чёрной дыры площади~$A$ достигается насыщение:
\begin{equation}
  \entropyS = \frac{\constBoltzmann A}{4\lengthPlanck^2}.
  \label{eq:bh-entropy}
\end{equation}
\begin{equation}
  \left[\frac{\constBoltzmann A}{4\lengthPlanck^2}\right]
  = \NotationUnitJoulePerKelvin
  = [\entropyS].
  \label{eq:bh-entropy-dim}
\end{equation}
Чтобы область той же шкалы не схлопнулась в горизонт, энергия удовлетворяет
\begin{equation}
  \energyE \le \frac{R\,\speedC^4}{2\constGrav},
  \qquad
  R_s=\frac{2\constGrav\,\energyE}{\speedC^4}.
  \label{eq:energy-bound}
\end{equation}
\end{proposition}

У шара радиуса~$R$ правая часть \eqref{eq:state-count} пропорциональна~$R^2$.
Набор независимых значений поля в каждой точке объёма рос бы как объём, то есть как~$R^3$,
а непрерывное значение в каждой точке $x\in\spaceRthree$ даёт несчётное множество конфигураций.
Общая теория относительности и квантовая теория поля локально друг другу не противоречат
и счёта \eqref{eq:state-count} не содержат.

Там, где одновременно существенны $\planckHbar$, $\constGrav$ и~$\speedC$, масштаб есть~$\lengthPlanck$ из \eqref{eq:bh-entropy}.
Поле в каждой точке счёт \eqref{eq:state-count} нарушает.
Следующий этаж --- запись того же сохранения и той же обратимости,
при которой спектр атома и два знака спина уже получаются.

